% !TeX root = ../../elegantbook-cn.tex


\section{切线}
导数的几何意义为函数图像在某点的切线斜率，因此我们可以借助导数求有关于函数图像的切线方程
\subsection{曲线$y=f(x)$在某点的切线方程}

设点$P(x_0,f(x_0))$在曲线$y=f(x)$上，且$f'(x_0)$存在，则曲线在点$P$处的切线斜率为
\[
k=f'(x_0).
\]
因此，切线方程为
\[
y-f(x_0)=f'(x_0)(x-x_0).
\]

求解时，先求导数$f'(x)$，再计算$f'(x_0)$和$f(x_0)$，最后代入点斜式方程即可。

\subsection{过某点与曲线$y=f(x)$相切的直线方程}

已知定点$P(x_0,y_0)$，求经过点$P$且与曲线$y=f(x)$相切的直线。由于点$P$不一定是切点，因此不能直接把$f'(x_0)$作为切线斜率。

设切点为$T(t,f(t))$，则切线斜率为$f'(t)$，切线方程为
\[
y-f(t)=f'(t)(x-t).
\]
因为切线经过点$P(x_0,y_0)$，所以
\[
y_0-f(t)=f'(t)(x_0-t).
\]
解出切点横坐标$t$后，将其代入
\[
y-f(t)=f'(t)(x-t)
\]
即可得到所求切线方程。方程可能有一个、多个或没有实数解，分别对应一条、多条或不存在满足条件的切线。



